
We will examine an example of an operator which will put our machinery
to work.

Our operator will be the \emph{momentum} operator
\begin{align}
  P_j \colon \Cal D_P &\to L^{2}(\R^{d}) \\
  \psi & \mapsto -\ic\hbar \partial_j\psi
\end{align}
For the purpose of showing the basic principles of the theory,
it is enough to take \(d=1\).

The first problem we must address is to choose the domain \(\Cal
D_P\). Of course, it cannot be the whole space \(L^{2}(\R)\).
Often, one is just ``given the operator'' without a domain, but really
one needs the domain to know whether the operator \(P\) is
self-adjoint. In fact, we want to pick a domain \(\Cal D_P\) such
that \(P\) \emph{is} self-adjoint, since it is supposed to be an
observable of a physical system.

In fact, instead of the real line \(\R\), we will first look at the
problem on a closed interval \([0,1]\), by choosing the domain to be
\begin{equation}
  \Cal D_P = \{\psi\in C^{1}[0,1] \mid \psi(0)=\psi(1)=0\}
,\end{equation}
where it will turns out that the operator cannot be made self-adjoint
in a unique way.
Then we will look at the same problem over the circle, by choosing the
domain
\begin{equation}
  \Cal D_P = \{\psi\in C^{1}[0,1] \mid \psi(0)=\psi(1)\}
,\end{equation}
where the momentum operator it will turn out to be essentially
self-adjoint, so it will have a unique self-adjoint extension.
In both cases, the Hilbert space under consideration will be \(\Cal
H=L^{2}[0,1]\).

\section{Absolutely continuous functions and Sobolev spaces}

We will encounter the spaces
\begin{equation}
  C^{1}[a,b] \subseteq H^{1}[a,b] \subseteq \AC[a,b]
.\end{equation}

\begin{definition}
  A function \(\psi\colon[a,b]\to\C\) is called \emph{absolutely
  continuous} if there is an integrable function
  \(\rho\colon[a,b]\to\C\) such that
  \begin{equation}
    \label{eq:abscont}
    \psi(x) = \psi(a) + \int_{a}^{x}\rho(y)\d y
  .\end{equation}
  The subspace of the elements \(\psi\in L^{2}[a,b]\) with an absolutely
  continuous representative is denoted \(\AC[a,b]\).
\end{definition}

Recall that integrable functions on \([a,b]\) are just \(L^{1}[a,b]\)
(after identifying almost-everywhere equal functions).

Note that, if \(\psi\in L^{1}[a,b]\) has an absolutely continuous
representative, then the integrable function \(\rho\) in equation
\eqref{eq:abscont} is unique up to almost-everywhere equality.
Hence, we have a differentiation map
\begin{equation}
  D \colon \AC[a,b]\to L^{1}[a,b]
\end{equation}
taking \(\psi\mapsto\rho\). We write \(\psi'=\rho\).

\begin{definition}
  The subspace \(H^{1}[a,b]\) of all \(\psi\in\AC[a,b]\) such that
  \(\psi'\in L^{2}[a,b]\) is called the \(1\)-th Sobolev space.
\end{definition}

\section{Momentum operator on a compact interval}

Consider the operator
\begin{align}
  P \colon \Cal D_P &\to L^{2}(\R^{d}) \\
  \psi & \mapsto -\ic \psi'
\end{align}
with domain
\begin{equation}
  \Cal D_P = \{\psi\in C^{1}[0,1] \mid \psi(0)=\psi(1)=0\}
.\end{equation}

The main question we want to tackle is:
\emph{is this operator self-adjoint?}

To even have a possibility of being self adjoint, it must be
symmetric. This is easier to check:
Take \(\psi,\phi\in\Cal D_P\). Using integration by parts, we have
\begin{align}
  \langle \psi,P\phi\rangle
  &= \int_{0}^{1} dx\ \bar\psi(x) (-\ic) \phi'(x) \\
  &= - \int_{0}^{1} dx \ \bar\psi'(x) (-\ic) \phi(x)
  + \cancelto{0}{\bar\psi(x)\phi(x) |_0^{1}} \label{eq:cancellation}\\
  &= \int_{0}^{1} dx \ \bar\psi'(x) \ic \phi(x) \\
  &= \int_{0}^{1} dx \ \overline{-\ic\psi'(x)} \phi(x) \\
  &= \langle P\psi,\phi\rangle
,\end{align}
so \(P\) is, indeed, symmetric.
We now calculate the adjoint \(P^{*}\).

\subsection{Calculating the adjoint on the interval}
Since \(P\) is symmetric, we know (\cref{prp:symext})
that \(P\subseteq P^{*}\), so \(P\) is self-adjoint iff the inclusion
\(\Cal D_P\subseteq \Cal D_{P^{*}}\) is an equality.
So, let's calculate the domain of \(P^*\).
An element \(\psi\in L^{2}[0,1]\) is in \(\Cal D_{P^{*}}\) iff
\begin{equation}
  \exists\eta\in L^{2}[0,1], \
  \forall\phi\in \Cal D_P, \
  \langle \psi, P\phi\rangle
  =
  \langle \eta, \phi\rangle,
\end{equation}
i.e.
\begin{equation}
  \exists\eta\in L^{2}[0,1], \
  \forall\phi\in \Cal D_P, \
  \int_{0}^{1}dx\, \bar\psi(x) (-\ic)\phi'(x)
  =
  \int_{0}^{1}dx\, \bar\eta(x) \phi(x).
\end{equation}
Now, every \(\eta\in L^{2}[0,1]\) is also integrable
(from measure theory),
so the function
\begin{equation}
  N(x) = \int_{0}^{x}\eta(x)\,dx
\end{equation}
is absolutely continuous with \(N'=\eta\). Hence,
\begin{align}
  \int_{0}^{1}dx\, \bar\eta(x) \phi(x)
  &= \int_{0}^{1}dx\, \bar N'(x) \phi(x) \\
  &= -\int_{0}^{1}dx\, \bar N(x) \phi'(x)
    + \cancelto{0}{\bar N(x)\phi(x)}
\end{align}
Hence, \(\psi\in L^{2}[0,1]\) is in \(\Cal D_{P^*}\) iff
\begin{equation}
  \exists N\in \AC[0,1], \
  \forall\phi\in \Cal D_P, \
  \int_{0}^{1}dx\, \bar\psi(x) (-\ic)\phi'(x)
  =
  -\int_{0}^{1}dx\, \bar N(x) \phi(x).
\end{equation}
i.e., iff
\begin{equation}
  \exists N\in \AC[0,1], \
  \forall\phi\in \Cal D_P, \
  \int_{0}^{1}dx\, \overline{(N(x)+\ic\psi(x))} \phi'(x)
  =0
\end{equation}
This condition means that there is an absolutely continuous
\(N\colon[0,1]\to\C\) such that \(N+\ic\psi \in \{\phi' \mid
\phi\in\Cal D_P\}^{\perp}=D(\Cal D_P)^{\perp}\).
Note that
\begin{equation}
  D(\Cal D_P) = \{f\in C^{0}[0,1] : \textstyle\int_{0}^{1}f(x)\d x=0\}
\end{equation}
and
\begin{equation}
  \overline{D(\Cal D_P)}
  = \{1\}^{\perp}
,\end{equation}
so \(D(\Cal D_P)^{\perp} = \langle 1\rangle\) are the constant
functions.
Hence, \(\psi\in\Cal D_{P^*}\) iff there is \(N\in\AC[0,1]\) such that
\(N+\ic\psi = c\) is a constant, but this happens iff \(\psi\in\AC[0,1]\).
Since we want \(P\psi=(-\ic)\psi'\in L^{2}[0,1]\), we conclude that
\(\Cal D_{P^*}=H^{1}[0,1]\).
We also know that the action of \(P^*\) is given by
\(P^{*}\psi=\eta=N'=(c-\ic\psi)'=-\ic\psi'\). In short, we calculated
that the adjoint of \(P\) is
\begin{align}
  P^{*}\colon H^{1}[0,1] &\to L^{2}[0,1] \\
  \psi &\mapsto -\ic\psi'
.\end{align}
Note that, as expected, \(\Cal D_P\subseteq \Cal D_{P^*}\), but this
inclusion is \emph{not} an equality, so \(P\) is \emph{not}
self-adjoint.

What hope do we have?
There is still the posibility that \(P\) is \emph{essentially self
adjoint}, i.e. that the closure \(P^{**}\) is self adjoint. This would
be the next best case, since that would mean that \(P^{**}\) is the
\emph{unique} self-adjoint extension of \(P\).

\subsection{Calculating the closure on the interval}

Recall that \(P\) is symmetric.
In this case we have
\begin{equation}
  P\subseteq P^{**}\subseteq P^{*}
\end{equation}
and we know that \(P^{**}\) is symmetric.
Note that the adjoint of \(P^{**}\) is \(P^{***}=P^{*}\)
(since \(P^{*}\) is already closed).
Hence, \(P^{**}\) is self adjoint iff the inclusion \(\Cal
D_{P^{**}}\subseteq \Cal D_{P^*}\) is an equality.
We now calculate the domain of \(P^{**}\).
An element \(\psi\in \Cal D_{P^*}=H^{1}[0,1]\) is in \(\Cal
D_{P^{**}}\) iff
\begin{equation}
  \exists\eta\in L^{2}[0,1], \
  \forall\phi\in\Cal D_{P^*}, \
  \langle \psi, P^{*}\phi \rangle
  =
  \langle \eta, \phi\rangle
.\end{equation}
Given such an element \(\psi\in\Cal D_{P^{**}}\), we have \(P^{**}\psi=\eta\).
Moreover, since \(P^{**}\subseteq P^*\), we also have
\(\eta=P^{**}\psi=P^*\psi\), so we obtain
\begin{equation}
  \forall\phi\in\Cal D_{P^*}, \
  \langle \psi, P^{*}\phi \rangle
  =
  \langle P^{*}\psi, \phi\rangle
.\end{equation}
In other words, if \(\psi\in\Cal D_{P^*}=H^{1}[0,1]\) is in \(\Cal
D_{P^{**}}\), then
\begin{equation}
  \forall\phi\in H^{1}[0,1], \
  \int_{0}^{1} dx\, \bar\psi(x)(-\ic)\phi'(x)
  =
  \int_{0}^{1} dx\, \overline{(-\ic)\psi'(x)}\phi(x)
.\end{equation}
Using integration by parts and rearranging, we see that both sides
differ by a boundary term, so this is equivalent to
\begin{equation}
  \forall\phi\in H^{1}[0,1], \
  \bar\psi(x)\phi(x)|_{0}^{1} = 0
\end{equation}
which happens iff \(\psi(0)=\psi(1)=0\).
Conversely, if \(\psi\in H^{1}[0,1]\) satisfies \(\psi(0)=\psi(1)=0\),
then \(\eta=-\ic\psi'\in L^{2}[0,1]\) satisfies, for all \(\phi\in
H^{1}[0,1]\)
\begin{align}
  \langle \psi, P^{*}\phi \rangle
  &= \int_{0}^{1} \bar\psi(x)(-\ic)\phi'(x)\d x \\
  &= -\int_{0}^{1} \bar\psi'(x)(-\ic)\phi(x)\d x
  + (-\ic)\bar\psi(x)\phi(x)|_{0}^{1} \\
  &= -\int_{0}^{1} \bar\psi'(x)(-\ic)\phi(x)\d x
  + (-\ic)\cancelto{0}{\bar\psi(x)\phi(x)|_{0}^{1}} \\
  &= \int_{0}^{1} \bar\psi'(x)\ic\phi(x)\d x \\
  &= \int_{0}^{1} \overline{-\ic\psi'(x)}\phi(x)\d x \\
  &= \langle \eta,\phi\rangle
,\end{align}
so \(\phi\in\Cal D_{P^{**}}\).
We conclude that
\begin{equation}
  \label{eq:domainpss}
  \Cal D_{P^{**}} = \{\psi\in H^{1}[0,1] : \psi(0)=\psi(1)=0\}
.\end{equation}
In particular, the inclusion \(\Cal D_{P^{**}}\subseteq \Cal
D_{P^*}\) is \emph{not} an equality.
(Note that, to reach this conclusion, we only
actually needed the inclusion \(\subseteq\) in equation
\eqref{eq:domainpss}, but we checked the other inclusion for fun).
Thus, we have shown that \(P^{**}\) is \emph{not} self adjoint, which
means that \(P\) is not even \emph{essentially} self-adjoint.

The situation looks bad.
Since \(P\) is not essentially self-adjoint, we cannot conclude that
it has a unique self-adjoint extension. However, it may be worse: it
could have no self-adjoint extension at all. To check for this
possibility, we calculate its defect indices.

\subsection{Calculating defect indices on the interval}

Recall that if an operator \(A\) has equal defect indices
\begin{equation}
  d_+ = \dim \ker(P^{*}-\ic),
  \qquad
  \qquad
  d_- = \dim \ker(P^{*}+\ic)
,\end{equation}
then \(A\) has a self-adjoint extension.
We calculate the kernels on the right hand sides, so we need to
determine the space of the \(\psi\in\Cal D_{P^*}=H^{1}[0,1]\) satisfying
\((P^*\mp\ic)\psi=0\), i.e. \(-\ic\psi'\mp\ic\psi=0\) or,
equivalently, \(\psi'\pm\psi=0\).
The solutions are all of the form \(\psi(x)=ae^{\mp x}\), so the
dimension of each kernel is \(1\). In particular, the defect indices
are equal, so there \emph{are} self-adjoint extensions of our operator \(P\).

In fact, it turns out that \(P\) has a family of self-adjoint
extensions, parametrized by \(U(1)\), but we won't determine it here.
Our current techniques are not enough.
We would need to learn about Cayley transforms.

An example where the operator \(P\) has \emph{no} self-adjoint
extensions is if we tried to do it on \([0,\infty)\) instead of
\([0,1]\). This means that, on a half line, there is no well-defined
momentum operator.

\section{Momentum operator on a circle}

Now we consider the operator
\begin{align}
  P \colon \Cal D_P &\to L^{2}[0,1] \\
  \psi &\mapsto -\ic\psi'
,\end{align}
where the domain is
\begin{equation}
  \Cal D_P = \{\psi\in C^{1}[0,1] \mid \psi(0)=\psi(1)\},
\end{equation}
without the condition that \(\psi\) attains the value \(0\) at the
boundary. Hence, this \(P\) is a proper extension
\begin{equation}
  \label{eq:extintcirc}
  P_{\text{interval}}
  \subsetneq
  P_{\text{circle}} = P
.\end{equation}

Again, we can check that \(P\) is symmetric. The calculation is the
same as for the interval, except that in step \eqref{eq:cancellation},
the boundary term \(\bar\psi(x)\phi(x)|_0^{1}\) vanishes for a
different reason: while for the interval every term was zero,
for the circle we only have \(\bar\psi(1)\psi(1)=\bar\psi(0)\phi(0)\).

\subsection{Calculating the adjoint on the circle}

Since \(P\) is symmetric, we know that \(P\subseteq P^{*}\).
Moreover, by equation \eqref{eq:extintcirc} it follows that
\(P^{*}\subseteq P^{*}_\text{interval}\). In particular, \(\Cal
D_{P^{*}}\subseteq H^{1}[0,1]\) and \(P^*\)
acts by \(P^*\psi=-\ic\psi'\).
Hence, taking \(\psi\in\Cal D_{P^{*}}\), we have
\begin{equation}
  \forall\phi\in\Cal D_P, \
  \langle P^*\psi, \phi \rangle
  =
  \langle \psi, P\phi \rangle
,\end{equation}
i.e.
\begin{equation}
  \forall\phi\in\Cal D_P, \
  \int_{0}^{1}dx\, \overline{-\ic\psi'(x)}\phi(x)
  =
  \int_{0}^{1}dx\, \overline{\psi(x)}(-\ic)\phi'(x)
.\end{equation}
Again, integrating by parts we see that both sides differ by a
boundary term \(\ic\bar\psi(x)\phi(x)|_{0}^{1}\), so their equality is
equivalent to the vanishing of the boundary term:
\begin{align}
  \forall\phi\in\Cal D_P, \
  \ic\bar\psi(x)\phi(x)|_{0}^{1} = 0
,\end{align}
and this happens iff \(\psi(0)=\psi(1)\).
Conversely, if \(\psi(0)=\psi(1)\), the element \(\eta=-\ic\psi'\)
satisfies the condition to be \(P^{*}\psi\), so \(\psi\in\Cal
D_{P^{*}}\).
Hence,
\begin{equation}
  \Cal D_{P^*} = \{\psi\in H^{1}[0,1] : \psi(0)=\psi(1)\}
.\end{equation}

In particular, the inclusion \(\Cal D_P\subseteq \Cal D_{P^{*}}\) is
not an equality, so \(P\) is \emph{not} self-adjoint.
Is it essentially self adjoint?

\subsection{Calculating the closure on the circle}

If \(P\) were essentially self-adjoint (i.e. if its closure \(P^{**}\)
were self-adjoint), then we would be done, since \(P^{**}\) would be
the unique self-adjoint extension of \(P\).
Again, since \(P^*\) is closed, the adjoint of \(P^{**}\) is
\(P^{***}=P^{*}\), so in view of the extensions
\begin{equation}
  P\subseteq P^{**} \subseteq P^*
\end{equation}
(which we have because \(P\) is symmetric),
it suffices to see whether the inclusion \(\Cal D_{P^{**}}\subseteq
\Cal D_{P^*}\) is an equality.
Given \(\psi\in\Cal D_{P^*}\) (i.e. \(\psi\in H^{1}[0,1]\) and
\(\psi(0)=\psi(1)\)), the function \(\eta=-\ic\psi'\in L^{2}[0,1]\)
satisfies,
\begin{align}
  \forall\phi\in\Cal D_{P^*}, \
  \langle \eta,\phi\rangle
  &= \int_{0}^{1}dx\, \overline{(-\ic)\psi'(x)}\phi(x) \\
  &= \ic\int_{0}^{1}dx\, \bar\psi'(x)\phi(x) \\
  &= -\ic\int_{0}^{1}dx\, \bar\psi(x)\phi'(x)
    + \cancelto{0}{\ic\bar\psi(x)\phi(x)|_{0}^{1}} \\
  &= \int_{0}^{1}dx\, \bar\psi(x)(-\ic)\phi'(x) \\
  &= \langle \psi,P^*\phi\rangle
,\end{align}
so \(\psi\in\Cal D_{P^{**}}\). This means that \(P^{**}=P^*\) so,
indeed, \(P^{**}\) is self-adjoint. It follows that it is the unique
self-adjoint extension of \(P\), and we are done!
Hence, the momentum operator on the circle is given by
\begin{align}
  \bar P \colon \{\psi\in H^1[0,1] : \psi(0)=\psi(1)\} &\to L^{2}[0,1]
  \\
  \psi &\mapsto -\ic\psi'
.\end{align}





